\section{Adaptive boundary summaries for repeated completion queries}
\label{sec:boundary-reuse}

Repeated observations of one suffix should share its topology.  The maintained
Euler backend now starts with direct closure walks and switches to a pair of
boundary path summaries when enough distinct queries reach the same stage.
A separate research implementation extends this idea to report 26's signed
Tait kernel.  Its common vertices are \emph{cut-face fragments}; this choice
makes both the quotient construction and its geometric checks depend on the
boundary rather than on another traversal of the entire suffix.
Neither improvement bounds the number of relative generators or proves a
general subexponential recognition algorithm.

\subsection{Two path summaries determine the raw Euler value}

Fix a suffix with $n$ crossings and $b$ boundary darts.  Its partial edge
involution is $\alpha$; at each crossing the opposite-dart pairing describes
link strands, and the zero-smoothing pairing describes the all-zero state.
Combining either local pairing with the internal edges gives disjoint closed
circles and paths whose endpoints lie on the boundary.  Store the number of
closed circles and the induced boundary matching for each of these two graphs.
Preparation takes $O(n+b)$ combinatorial work and temporary space, while the
retained summaries have size $O(b)$ apart from the integer counts.

For a proposed boundary matching $m$, the circles of the completed graph
are exactly the stored closed circles plus the alternating cycles of the
stored boundary matching with $m$.  Thus two boundary walks compute the link
component count $c$ and the all-zero circle count $c_0$.  Coverage is checked:
every boundary label must occur once, with no missing, repeated, or unknown
endpoint.  No suffix dart need be revisited by these queries.

Here is why these counts suffice even for the sign.  Choose any orientation
of the classical completion, let $h$ be the number of one-smoothings in its
oriented Seifert state, and let $s$ be its number of Seifert circles.  Changing
one smoothing changes circle parity, giving $s\equiv c_0+h\pmod2$.
The orientable Seifert surface has Euler characteristic $s-n$ and $c$
boundary components, hence $s-n\equiv c\pmod2$.  This also holds for a
disconnected surface.  Therefore
\[
 h\equiv n+c+c_0\pmod2,\qquad
 \chi_{\mathrm{raw,unred}}(m)=(-1)^{n+c+c_0}2^c.
\]
This is the unnormalized Euler value needed by the existing relative scan.
The production caller supplies classical completions from a valid scan;
these two path summaries alone are not a test for virtual planarity.
At the empty scalar terminal the convention is $n=c=c_0=0$, with value one;
it does not denote an implicitly added crossingless unknot.

Each subsequent query uses $O(b)$ combinatorial operations.  This is not
an $O(b)$ bound on bit complexity independent of $n$: the integer $2^c$ itself
has $c+1=O(n+b)$ bits.  The implementation retains exact integers and the
existing coefficient-size statistics.

\subsection{When the maintained backend switches}

The summaries are implemented in
\path{fast/fastunknot/boundary_connectivity.py}.

\code{AdaptiveClosureEuler} in \path{fast/fastunknot/euler_scan.py} uses the
independent orientation-based \code{ClosureEuler} for the first four distinct,
uncached queries at a stage.  Later queries prepare the summaries if $b\le n$;
a wide boundary on a short suffix keeps direct evaluation.  This is an
empirical setup policy, not an optimal stopping theorem.  A low-level
\code{compression\_after=0} option forces eager preparation for diagnostics.
The public Euler backend uses the adaptive policy; the standard recognizer
backend remains unchanged.

Changing the stage discards the active summaries.  Cached exact values do
not consume another query allowance.  The global deadline is checked even
on a cache hit, and the local query cap is checked before preparation.
A summary is published only after both path traversals finish, so an
interrupted setup cannot leave a usable half-summary.  Local inference
exhaustion continues the same exact saturated scan, while global exhaustion
returns \code{UNKNOWN}, as before.  The direct closure evaluator and the
older suffix-smoothing recurrence remain independent references.

Five new tests cover random classical completions and their mirrors,
independent recurrence values and explicit circle counts, switching and
stage changes, the wide-boundary guard, cache and budget behavior, malformed
coverage, and interruption after the first of the two path summaries.
The full maintained suite passes all 572 tests; logs are retained in
\path{data/boundary-connectivity-tests.txt} and
\path{data/boundary-connectivity-integrated-tests.txt}.

\subsection{A common Tait graph made from cut-face fragments}

The research implementation \path{fast/determinant_research/boundary_tait.py}
starts with a checkerboard coloring of the original classical diagram.
Write $\rho$ for cyclic rotation of a crossing's darts.  A completed face
follows the permutation $\rho\alpha$.  Retain only its edges arising from
internal suffix edges.  Each resulting component is either a closed face
cycle or an open face path.  These \emph{fragments}, rather than entire
original regions, are the common graph vertices.  In particular, a region
that splits when the prefix is removed is represented by separate fragments.

For a boundary label $a$ with dart $d_a$, store the fragment containing
$d_a$ as its outgoing endpoint and the fragment containing $\rho(d_a)$ as
its incoming endpoint.  Pairing labels $a,b$ inserts exactly the two face
connections
\[
 d_a\longrightarrow\rho(d_b),\qquad
 d_b\longrightarrow\rho(d_a).
\]
Consequently every completed face is obtained by identifying these open
fragments, while every closed fragment remains untouched.  This proves the
face reconstruction directly from the partial permutation; it does not
assume that the original region partition survives a new closure.

There are at most $2b$ open fragments.  A query runs union-find only on them
and on the suffix projection components incident to the boundary.  Counts
of closed fragments and closed projection components are stored once.
Every inserted face connection must preserve the inherited checkerboard
color.  If $F$ is the resulting face count and $k$ the resulting projection
component count, the query also requires
\[
 F=n+2k.
\]
The completed rotation system has $V=n$ and $E=2n$.  Its orientable ribbon
components have total Euler characteristic $F-n=2k-2\sum g_i$; the displayed
check is therefore equivalent to every $g_i=0$.  It certifies spherical
completion for accepted queries.  Color incompatibility is a decline of
this particular representation, even when a completion is classical;
it is never evidence that a knot is nontrivial.

Choose the smaller shading by fragment count.  Each crossing contributes
one signed edge between its two black fragments, with the same local weight
as in Section~\ref{sec:determinant-continuations}.  The black open fragments
are terminals.  All identifications of the common graph occur among these
terminals; black closed fragments are fixed interior vertices.  Thus an
accepted connected completion has exactly the quotient Tait graph required
by report 26's terminal kernel.  Its phase is recomputed from the resulting
black vertex count $v$:
\[
 J_{\mathrm{raw}}(i)=(-i)^{B+v-1}\tau(G/m).
\]
The crossing sum $B$ includes loop edges even though loops contribute zero
to the Laplacian.  If the completed projection is disconnected, the
specialization is zero and no kernel is needed.  If there is no boundary,
one black vertex serves as the artificial grounding terminal.

The common Laplacian is prepared once; the delivered exact rational kernel
is built lazily.  Canonical terminal partitions cache completed integer
cofactors, so distinct matchings inducing the same partition share a
value.  Phase and topology counts are still computed for each matching.
The two boundary path summaries give the raw value at one and the parity
$\epsilon=c_0-1\pmod2$.  Together with the value at $i$, they reconstruct
the four marked residues by the existing $(a+d)/2,(a-d)/2$ formula.

The geometric query costs $O(b\,\alpha(b))$ combinatorial operations.
For the determinant, the singular-safe elimination described earlier leaves
a border of dimension $r+q$, where $r$ is interior nullity and $q$ is one
less than the number of terminal blocks.  If $r>q$, the cofactor vanishes;
otherwise $r+q=O(b)$.  This controls matrix dimension, not coefficient bit
length or the number of scan queries.  Rational arithmetic can still depend
polynomially on the original suffix size.  Production integration also
requires an interruptible maintained kernel and useful amortization on
actual observer traces.  The present prototype imports report 26's kernel
and has no production deadline or local-work interface; production shadow
observations continue to compute their closures directly.

\subsection{Independent geometry audit}

\path{fast/determinant_research/audit_boundary_tait.py} repeats all 4,503
actual completion queries on 582 proper-prefix stages of 78 retained input
diagrams.  Every accepted value agrees with an independently rebuilt
completed diagram, including component counts, zero circles, raw Euler
values, and the Gaussian value at $i$.  There are 1,317 disconnected
projections, 15 singular stage kernels, maximum terminal count ten, and
maximum interior nullity two.  Kernel certificates are checked independently.

A further exhaustive pairing audit covers boundary sizes four through eight
in the first 20 retained diagrams: among 3,468 pairings, 745 are accepted,
472 classical completions are declined because the inherited coloring does
not fit, and 2,251 nonclassical completions are rejected.  Every accepted
pairing passes the independent classical completion check and both value
comparisons.  The scope, exact inputs, and source hashes are recorded in
\code{data/boundary-tait-audit.json}.  These finite checks support the
implementation; the fragment and Euler-characteristic arguments above
supply the general accepted-query contract.

\subsection{Measured reuse and its limits}

\path{fast/benchmark_boundary.py} uses seven shuffled paired rounds with
identical control arms and one excluded warmup.  Every timed sample includes
fresh engine setup.  The isolated workload consists of 200 distinct actual
matchings from one nine-crossing prefix, all with 18 boundary darts.  Appended
cancelling braid pairs increase the suffix from 5 to 37, 133, and 517
crossings without changing the knot.  Matching-stream discovery is outside
the primitive timings.  These deliberately easy added crossings isolate
reuse; they are not a hard family for complete recognition.

\input{tables/boundary_streams.tex}

The adaptive Euler policy achieves up to $15.8$ times speedup, including its
first four direct queries and summary setup.  Eager setup reaches $23.6$
times on the longest suffix.  On the shortest suffix, the guard prevents
preparation but the adaptive wrapper still costs about 15 percent in this
isolated 200-query workload.  This overhead is retained in the table.

The research determinant path reaches $53.9$ times speedup, including
coloring, fragment geometry, kernel construction, and partition-cache setup.
However, all four of these common graphs have nine vertices, all nine are
terminals, and interior nullity is zero.  These gains therefore demonstrate
boundary geometry and partition reuse, not amortization of a nontrivial
interior elimination or a singular-kernel runtime advantage.  Only 61 or
96 distinct cofactors are computed for each 200-matching stream; disconnected
queries also avoid determinants.  Independent kernel-certificate checking
is performed in the separate audit, outside these timings.

\input{tables/boundary_pipeline.tex}

For complete raw Euler scans, the baseline loads the earlier
\code{euler\_scan.py} from revision \code{f0179ab50}; the other dependencies
are shared.  Status, method, and observation stage agree in every paired
run.  None of these seven raw scans prepares a boundary summary.  They show
no consistent gain, and the stress case is about four percent slower.
For complete recognition with the Euler backend requested, earlier filters
decide every example, so no Euler evaluator runs.  No whole-recognizer
speedup is established.  The raw samples, per-arm evidence, inputs, and
source hashes are retained in
\path{fast/results/boundary_reuse_20261008.json}.  Future scheduling decisions
must distinguish frequent cheap boundary queries from early-stopping scans;
neither a favorable microbenchmark nor an adaptive switch supplies the
missing universal width bound.
